% META: {"id": "TSPE-PROBA-EV-A-corrige", "chapitre": "TSPE-PROBABILITES", "type_objet": "corrige_evaluation", "evaluation_ref": "TSPE-PROBA-EV-A", "version": "A", "capacites_codes": ["C1", "C2", "C4", "C5", "C9"], "sources_inspiration": [], "status": "approved"}
% BEGIN-VERIFY
% from sympy import Rational, binomial
% p1 = Rational(5,8)*Rational(3,7) + Rational(3,8)*Rational(5,7)
% assert p1 == Rational(15, 28)
% p2 = binomial(4,1)*Rational(1,4)**1*Rational(3,4)**3
% assert p2 == Rational(27, 64)
%
% from sympy import N
% # Schema de Bernoulli : 20 epreuves independantes, succes de probabilite 1/4.
% n, p = 20, Rational(1, 4)
% def P_egal(k):
%     return binomial(n, k) * p**k * (1 - p)**(n - k)
% assert sum(P_egal(k) for k in range(n + 1)) == 1
% assert abs(N(P_egal(5)) - 0.2023) < 0.0001
% cumul = sum(P_egal(k) for k in range(6))
% assert abs(N(cumul) - 0.6172) < 0.0001
% # Plus petit seuil a 95 %.
% seuil = next(k for k in range(n + 1)
%              if sum(P_egal(j) for j in range(k + 1)) >= Rational(95, 100))
% assert seuil == 8
% assert sum(P_egal(j) for j in range(seuil)) < Rational(95, 100)
% END-VERIFY
\begin{corrige}{TSPE-PROBA-EV-A-EX1}
\textbf{Q1.} Premier tirage : $P(R)=\dfrac{5}{8}$, $P(B)=\dfrac{3}{8}$. Deuxième tirage (sans remise) : si $R$ au premier tirage, $P_R(B)=\dfrac{3}{7}$ ; si $B$, $P_B(R)=\dfrac{5}{7}$.

\textbf{Q2.} $P(\text{couleurs différentes}) = P(R_1)\times P_{R_1}(B_2) + P(B_1)\times P_{B_1}(R_2) = \dfrac{5}{8}\times\dfrac{3}{7}+\dfrac{3}{8}\times\dfrac{5}{7} = \dfrac{15}{56}+\dfrac{15}{56} = \boxed{\dfrac{15}{28}}$.
\end{corrige}

\begin{corrige}{TSPE-PROBA-EV-A-EX2}
Les chemins avec exactement $1$ succès parmi $4$ epreuves sont au nombre de $\dbinom{4}{1}=4$ (SEEE, ESEE, EESE, EEES), chacun de probabilité $p^1(1-p)^3 = 0{,}25\times(0{,}75)^3$.

\[
  P(X=1) = \binom{4}{1}(0{,}25)^1(0{,}75)^3 = 4\times0{,}25\times0{,}421875 = \boxed{\dfrac{27}{64}} \approx 0{,}4219.
\]
\end{corrige}

\begin{corrige}{TSPE-PROBA-EV-A-EX3}
$E(X) = np = 25\times0{,}6 = \boxed{15}$. $V(X) = np(1-p) = 25\times0{,}6\times0{,}4 = \boxed{6}$.
\end{corrige}

\begin{corrige}{TSPE-PROBA-EV-A-EX4}
\textbf{Q1.} L'épreuve est « répondre à une question » ; le succès est « la réponse est correcte », de probabilité $p=\dfrac14$ puisque le candidat choisit au hasard parmi quatre réponses dont une seule convient. Les questions sont \textbf{indépendantes} et de même probabilité de succès : c'est un schéma de Bernoulli à $n=20$ épreuves, donc $X\sim\mathcal{B}\left(20\,;\dfrac14\right)$.

\textbf{Q2.} $P(X=5)=\dbinom{20}{5}\left(\dfrac14\right)^5\left(\dfrac34\right)^{15}\approx0{,}2023$, et
$P(X\leq5)=\displaystyle\sum_{k=0}^{5}P(X=k)\approx0{,}6172$.

\textbf{Q3.} En cumulant, $P(X\leq7)\approx0{,}898<0{,}95$ tandis que $P(X\leq8)\approx0{,}959\geq0{,}95$ : le plus petit entier est $k=8$. Autrement dit, un candidat répondant entièrement au hasard obtient au plus $8$ bonnes réponses sur $20$ dans $95\%$ des cas : dépasser ce seuil est peu compatible avec le seul hasard.
\end{corrige}
